% =============================================================================
% Experiments section for the multi-agent GRPO paper (NeurIPS style).
%
% Drop into your main .tex with
%   % =============================================================================
% Experiments section for the multi-agent GRPO paper (NeurIPS style).
%
% Drop into your main .tex with
%   % =============================================================================
% Experiments section for the multi-agent GRPO paper (NeurIPS style).
%
% Drop into your main .tex with
%   % =============================================================================
% Experiments section for the multi-agent GRPO paper (NeurIPS style).
%
% Drop into your main .tex with
%   \input{paper_experiments_section}
% and ensure the following are also \input or \include'd before this file:
%   - paper_results_filled.tex   (Tables 1, 2 and 3)
%   - the figure block for       paper_figures/case_improvement_bars.pdf
%                                paper_figures/scaling_with_lines.pdf
%
% Required packages (NeurIPS standard, already used by the tables file):
%   \usepackage{booktabs}
%   \usepackage{multirow}
%   \usepackage{makecell}
% =============================================================================

\section{Experiments}
\label{sec:experiments}

We evaluate DPA-GRPO on TaxCalcBench TY24~\citep{bock2025taxcalcbench}, a
line-level federal-1040 reasoning benchmark, and compare against zero-shot
and trained single-agent reinforcement-learning baselines at two backbone
sizes (Qwen3-4B and Qwen3-8B). Our experiments are designed to answer
four questions, one per subsection: \emph{(i)} how is the two-player
problem framed and trained on this benchmark
(\S\ref{sec:exp-setup})?\ \emph{(ii)} how does DPA-GRPO perform overall,
and how does it scale with backbone size and training-case diversity
(\S\ref{sec:exp-results})?\ \emph{(iii)} is the gain attributable to the
two-player credit assignment, or could a stronger single-agent RL recipe
close it (\S\ref{sec:exp-baselines})?\ and \emph{(iv)} which line-level
dynamics actually change during training, and do they match the case
structure that DPA-GRPO is designed to exploit
(\S\ref{sec:exp-cases})?

\subsection{Problem setup}
\label{sec:exp-setup}

\paragraph{Benchmark.}
TaxCalcBench TY24~\citep{bock2025taxcalcbench} is a held-out collection
of U.S.\ federal tax returns whose ground truth was produced by a
deterministic tax engine. Each case provides the W-2/1099/Schedule
inputs and a 19-line gold 1040 (lines 1a, 9, 10, 11, 12, 15, 16, 19, 24,
25d, 26, 27, 28, 29, 32, 33, 34, 35a, 37). We use the official
train/test split with \texttt{split\_seed}=42 and
\texttt{test\_fraction}=0.2, yielding $n{=}8$ held-out test cases and a
total of $8\times 19 = 152$ independent line-level decisions per
evaluation. All accuracies in this section are computed against the
exact engine output for these 152 lines.

\paragraph{Two-player loop.}
Following Section~\ref{sec:method}, every line of the 1040 is treated as
a separate two-player game. The \emph{Generator} produces a draft value
for the line; the \emph{Verifier} either remains silent or raises a
\emph{Self-Assessment Critique} (SAC). When an SAC is raised, the
Generator emits a revised value and a binary \textsc{keep}/\textsc{revise}
decision. Both agents are LoRA adapters over the same frozen backbone, so
the only parameters trained by reinforcement learning are the two
adapters. The eight terminal states of the game collapse into the
mutually-exclusive case taxonomy used throughout the paper:
\textbf{Case~1} (Generator correct, Verifier silent),
\textbf{Cases~2--3} (Verifier intervenes when the original draft is
already correct),
\textbf{Case~4} (\emph{missed bug}: Generator wrong, Verifier silent),
\textbf{Cases~5a/5b} (Verifier intervenes on a wrong draft and the
revise either lands or misses gold),
and \textbf{Cases~6a/6b} (Generator overrides the Verifier; 6a is a
correct override, 6b an incorrect one).

\paragraph{Training.}
Both adapters are updated with 2-GRPO: one GRPO update on the Generator,
then one on the Verifier, repeated. Each update samples a group of
$G{=}2$ trajectories per case at temperature $0.2$ and clips the
importance ratio at $\varepsilon{=}0.2$; the shared reward is the
case-conditional shaping described in Section~\ref{sec:method}, coupled
with KL regularisation against a frozen reference
($\beta_{\text{KL}}{=}0.04$). We sweep \emph{training cases per
gradient step} $\in\{3, 5, 10\}$ for both backbones; all other
hyperparameters are held fixed (LoRA rank 16, learning rate
$1{\times}10^{-5}$, gradient checkpointing, fp16, single H100). Training
runs for at most 30 GRPO steps, with early stopping on held-out
post-revise accuracy (patience 2, minimum 10 steps).

\paragraph{Metrics.}
We report two accuracy metrics over the 152 line decisions on the
held-out test split. \textbf{Strict accuracy} is the fraction on which
the Generator's \emph{first} draft matches the gold engine output;
\textbf{post-revise accuracy} scores the Generator's \emph{final} value
after the optional revise. For zero-shot and trained single-agent
baselines the Generator is the only agent, so strict and post-revise
coincide; for DPA-GRPO we report post-revise accuracy unless otherwise
stated, since this is the metric the two-player loop optimises. We
additionally track \emph{decision accuracy} (the fraction of SAC-flagged
lines on which the Generator's keep/revise decision matches the
oracle correct-direction move) and the \emph{case-histogram} fractions
that decompose the 152 decisions into the taxonomy above. To control for
training-data leakage we also tabulate per-case training-time accuracy
on three representative training cases, \textsc{Train-A}/\textsc{B}/\textsc{C}
(see Table~\ref{tab:main-results} for full identifiers), encountered
by every trained run because they share \texttt{split\_seed}=42.

\subsection{Experimental results}
\label{sec:exp-results}

Table~\ref{tab:main-results} reports test post-revise accuracy and
training-time accuracy on \textsc{Train-A}/\textsc{B}/\textsc{C} for
every method we trained or evaluated. Three findings stand out.

\paragraph{DPA-GRPO is the best method at both scales.}
At Qwen3-4B, DPA-GRPO reaches $0.533$ test post-revise accuracy, which
already \emph{matches} the zero-shot accuracy of the next backbone size
(Qwen3-8B zero-shot, $0.533$). At Qwen3-8B, DPA-GRPO reaches $0.572$,
the highest entry in the table; this is a $+3.9$\,pp absolute
improvement over Qwen3-8B zero-shot ($0.533$) and a $+1.9$\,pp
improvement over the strongest single-agent RL baseline (generator-only
DAPO 8B, $0.553$). The same ranking holds on every per-case
training-time metric we tabulate: at 8B, DPA-GRPO is best on
\textsc{Train-A} ($0.895$ vs.\ $0.842$), \textsc{Train-B} ($0.684$ vs.\
$0.632$), and \textsc{Train-C} ($0.579$ vs.\ $0.526$), each a
$+5{-}6$\,pp improvement over the strongest single-agent 8B baseline.
The best DPA-GRPO configuration differs by backbone (4B prefers $10$
cases per step; 8B prefers $3$); the full sweep is reported in
Table~\ref{tab:ours-ablation}, and we adopt those bolded rows as the
headline numbers throughout the paper.

\paragraph{Scaling with training data.}
Figure~\ref{fig:scaling} plots accuracy against the cumulative number of
training-line decisions ingested during training, so that runs with
different \emph{cases per step} settings are placed on a common
$x$-axis. DPA-GRPO surpasses every untrained baseline once roughly
$2{\times}10^{3}$ line decisions have been seen, and continues to
improve when given more data; generator-only GRPO (the strongest
single-agent baseline at 4B) plateaus earlier and at a lower ceiling.
The gap between DPA-GRPO and generator-only GRPO is therefore not a
data-quantity artefact of unequal training budgets, but a property of
the objective itself.

\paragraph{Scaling with backbone size.}
The two backbones tell complementary stories. At 4B, the \textsc{Train-A}
case is far from saturation ($0.526$ zero-shot), so most of the lift
DPA-GRPO obtains there ($+21.1$\,pp to $0.737$) is genuine learning.
At 8B, \textsc{Train-A} is already near-saturating zero-shot ($0.842$),
so the marginal lift on that case is small ($+5.3$\,pp to $0.895$);
the larger backbone's headroom appears instead on \textsc{Train-B} and
\textsc{Train-C}, the two harder cases, where DPA-GRPO 8B delivers
$+5.2$ and $+5.3$\,pp over the strongest single-agent 8B baseline.
This is the pattern one would hope for: as the Generator gets stronger,
DPA-GRPO's gains migrate to harder line types rather than disappearing.

\subsection{Comparisons with baselines}
\label{sec:exp-baselines}

We compare DPA-GRPO against three families of baselines, all using the
same Qwen3-4B and Qwen3-8B backbones, the same train/test split, and the
same LoRA hyperparameters. The first family is untrained; the second is
\emph{generator-only}---a single LoRA adapter trained to maximise strict
accuracy on the gold draft, with no Verifier; and the third is the
trained, two-player DPA-GRPO.

\paragraph{Untrained baselines.}
The zero-shot Generator already produces calibrated drafts on this
benchmark, with strict accuracy $0.414$ at 4B and $0.533$ at 8B
(Table~\ref{tab:main-results}, top block). Because the Verifier loop is
part of our trained method only, the zero-shot rows are intentionally
single-agent: any post-revise accuracy attributed to a zero-shot row
would conflate the loop with the underlying model. This makes
DPA-GRPO's $+11.9$\,pp lift at 4B and $+3.9$\,pp lift at 8B a
conservative measure of the parameter-update contribution.

\paragraph{Generator-only RL baselines (GRPO, DAPO, GSPO).}
We re-implemented three single-agent RL baselines that share the GRPO
trajectory-grouping recipe but differ in their objective: \emph{vanilla}
GRPO~\citep{shao2024deepseekmath} with the same clip-and-KL setup as
DPA-GRPO's Generator branch; \emph{DAPO}~\citep{yu2025dapo}, which
decouples the policy clip into asymmetric
$\varepsilon_{\text{low}}$/$\varepsilon_{\text{high}}$ bounds and adds
dynamic sampling; and \emph{GSPO}~\citep{zheng2025gspo}, which replaces
the per-token policy ratio with a per-sequence one. All three are
trained against the same gold-draft outcome reward as DPA-GRPO's
Generator branch and are matched on wall-clock budget, so the
comparison isolates the contribution of the two-player credit
assignment. At 4B, DPA-GRPO ($0.533$) beats vanilla GRPO ($0.520$),
DAPO ($0.467$), and GSPO ($0.461$). At 8B, DPA-GRPO ($0.572$) beats
vanilla GRPO ($0.539$) and DAPO ($0.553$); the GSPO 8B run is in
progress and will be reported in the camera-ready.

\paragraph{The two-player loop also stabilises the ranking.}
Beyond the absolute numbers, the ranking among single-agent baselines
\emph{flips} between scales: DAPO is the worst single-agent recipe at
4B ($0.467$) but the best at 8B ($0.553$), and GRPO does the
reverse. DPA-GRPO is the strict winner at both scales and on every
\textsc{Train-A/B/C} column we tabulate. We read this as evidence that
the two-player credit assignment is not interchangeable with a sharper
single-agent objective; it provides a complementary stabilising signal
that single-agent recipes cannot match by tuning alone.

\paragraph{Why a Verifier helps.}
A generator-only learner only ever moves probability mass between
``correct'' and ``incorrect'' draft outcomes; it has no signal that
distinguishes a \emph{confidently} correct line from a fragile one. The
two-player loop adds a Verifier whose pass-through behaviour is itself
trained, so a Generator line that is correct \emph{but} would be
overturned by an unconfident Verifier is no longer interchangeable with
a confidently-correct one in the joint reward. We make this mechanism
concrete next.

\subsection{Case discovery: which line-level dynamics change?}
\label{sec:exp-cases}

A natural worry with the headline numbers in
Table~\ref{tab:main-results} is that the post-revise gain comes from a
single, narrow mechanism: the Verifier rescuing already-wrong drafts at
test time. If that were the only thing happening, DPA-GRPO would amount
to a learned form of test-time refinement and should be compared
against post-hoc self-revision rather than against single-agent RL. To
rule this out we decompose every evaluated line into the eight-case
taxonomy of \S\ref{sec:exp-setup} and ask which buckets actually move
during training.

Figure~\ref{fig:case-improvement} reports the case-mix shift
\emph{during training}: for our two productive 4B runs, we compare the
first-5-step rolling window (baseline; the joint policy before
DPA-GRPO updates have moved either head appreciably) against the most
favourable post-update 5-step window on the same training-rollout
stream. Both endpoints share the same Verifier-firing regime, so the
comparison isolates what the joint policy \emph{learns} during training
rather than how often the Verifier intervenes. Consequently, the
percentages here are training-time within-run shifts and differ from
the test-set case mix in Table~\ref{tab:main-results} (single held-out
checkpoint); the two answer complementary questions
(\emph{what does the 2-GRPO loop learn to do?} vs.\ \emph{how does the
trained loop behave on held-out data?}). Three movements emerge, all
in the \emph{intended} direction:

\begin{itemize}
\item \textbf{Case~1 (Generator correct, Verifier silent) rises by
  $+15.1\%$.} Case~1 is the desired terminal state---the Generator
  is right and the Verifier correctly does nothing. The lift confirms
  that the Verifier's joint pressure on the Generator is producing
  more confidently-correct drafts, not just more revisions.
\item \textbf{Case~4 (\emph{missed bug}) drops by $-18.1\%$.} Case~4
  is the dominant failure mode of the initial joint policy and the only
  bucket where neither agent can be retroactively credited with a save.
  The drop here is the largest single movement we observe and is the
  mechanism by which DPA-GRPO improves \emph{strict} accuracy, not just
  post-revise accuracy.
\item \textbf{Case~6a (Generator skips a bad revision) rises by
  $+7.4\%$.} In Case~6a the draft is wrong and the Verifier correctly
  raises SAC, but the Generator's own candidate revision would also be
  wrong; the Generator chooses KEEP rather than swapping in a
  no-better revision. This is the only bucket that explicitly
  measures the Generator's \emph{discriminative} skill against its
  own revision proposals, and the rise rules out the trivial
  alternative of an always-revise Generator that would post-hoc trade
  strict accuracy for post-revise accuracy.
\end{itemize}

The three movements are mutually reinforcing rather than redundant. The
Generator becomes more right (Case~1, Case~4) and also better
calibrated against its own revision proposals (Case~6a), while the
Verifier becomes more accurate at flagging genuine bugs (the only way
Case~4 can drop
this far). The same direction of movement holds on the 8B runs
(Appendix~\ref{app:case-trends}); the within-run magnitudes are smaller
because the 8B early-training baseline already sits in a
less-Case-4-heavy regime, so there is simply less probability mass left
to redistribute during the productive training window. Taken with the test
accuracy gains in Table~\ref{tab:main-results}, this answers
question~(iv): the headline accuracy improvement is reflected in real,
predicted shifts in the line-level case mix, not in a single
narrow refinement mechanism.

% and ensure the following are also \input or \include'd before this file:
%   - paper_results_filled.tex   (Tables 1, 2 and 3)
%   - the figure block for       paper_figures/case_improvement_bars.pdf
%                                paper_figures/scaling_with_lines.pdf
%
% Required packages (NeurIPS standard, already used by the tables file):
%   \usepackage{booktabs}
%   \usepackage{multirow}
%   \usepackage{makecell}
% =============================================================================

\section{Experiments}
\label{sec:experiments}

We evaluate DPA-GRPO on TaxCalcBench TY24~\citep{bock2025taxcalcbench}, a
line-level federal-1040 reasoning benchmark, and compare against zero-shot
and trained single-agent reinforcement-learning baselines at two backbone
sizes (Qwen3-4B and Qwen3-8B). Our experiments are designed to answer
four questions, one per subsection: \emph{(i)} how is the two-player
problem framed and trained on this benchmark
(\S\ref{sec:exp-setup})?\ \emph{(ii)} how does DPA-GRPO perform overall,
and how does it scale with backbone size and training-case diversity
(\S\ref{sec:exp-results})?\ \emph{(iii)} is the gain attributable to the
two-player credit assignment, or could a stronger single-agent RL recipe
close it (\S\ref{sec:exp-baselines})?\ and \emph{(iv)} which line-level
dynamics actually change during training, and do they match the case
structure that DPA-GRPO is designed to exploit
(\S\ref{sec:exp-cases})?

\subsection{Problem setup}
\label{sec:exp-setup}

\paragraph{Benchmark.}
TaxCalcBench TY24~\citep{bock2025taxcalcbench} is a held-out collection
of U.S.\ federal tax returns whose ground truth was produced by a
deterministic tax engine. Each case provides the W-2/1099/Schedule
inputs and a 19-line gold 1040 (lines 1a, 9, 10, 11, 12, 15, 16, 19, 24,
25d, 26, 27, 28, 29, 32, 33, 34, 35a, 37). We use the official
train/test split with \texttt{split\_seed}=42 and
\texttt{test\_fraction}=0.2, yielding $n{=}8$ held-out test cases and a
total of $8\times 19 = 152$ independent line-level decisions per
evaluation. All accuracies in this section are computed against the
exact engine output for these 152 lines.

\paragraph{Two-player loop.}
Following Section~\ref{sec:method}, every line of the 1040 is treated as
a separate two-player game. The \emph{Generator} produces a draft value
for the line; the \emph{Verifier} either remains silent or raises a
\emph{Self-Assessment Critique} (SAC). When an SAC is raised, the
Generator emits a revised value and a binary \textsc{keep}/\textsc{revise}
decision. Both agents are LoRA adapters over the same frozen backbone, so
the only parameters trained by reinforcement learning are the two
adapters. The eight terminal states of the game collapse into the
mutually-exclusive case taxonomy used throughout the paper:
\textbf{Case~1} (Generator correct, Verifier silent),
\textbf{Cases~2--3} (Verifier intervenes when the original draft is
already correct),
\textbf{Case~4} (\emph{missed bug}: Generator wrong, Verifier silent),
\textbf{Cases~5a/5b} (Verifier intervenes on a wrong draft and the
revise either lands or misses gold),
and \textbf{Cases~6a/6b} (Generator overrides the Verifier; 6a is a
correct override, 6b an incorrect one).

\paragraph{Training.}
Both adapters are updated with 2-GRPO: one GRPO update on the Generator,
then one on the Verifier, repeated. Each update samples a group of
$G{=}2$ trajectories per case at temperature $0.2$ and clips the
importance ratio at $\varepsilon{=}0.2$; the shared reward is the
case-conditional shaping described in Section~\ref{sec:method}, coupled
with KL regularisation against a frozen reference
($\beta_{\text{KL}}{=}0.04$). We sweep \emph{training cases per
gradient step} $\in\{3, 5, 10\}$ for both backbones; all other
hyperparameters are held fixed (LoRA rank 16, learning rate
$1{\times}10^{-5}$, gradient checkpointing, fp16, single H100). Training
runs for at most 30 GRPO steps, with early stopping on held-out
post-revise accuracy (patience 2, minimum 10 steps).

\paragraph{Metrics.}
We report two accuracy metrics over the 152 line decisions on the
held-out test split. \textbf{Strict accuracy} is the fraction on which
the Generator's \emph{first} draft matches the gold engine output;
\textbf{post-revise accuracy} scores the Generator's \emph{final} value
after the optional revise. For zero-shot and trained single-agent
baselines the Generator is the only agent, so strict and post-revise
coincide; for DPA-GRPO we report post-revise accuracy unless otherwise
stated, since this is the metric the two-player loop optimises. We
additionally track \emph{decision accuracy} (the fraction of SAC-flagged
lines on which the Generator's keep/revise decision matches the
oracle correct-direction move) and the \emph{case-histogram} fractions
that decompose the 152 decisions into the taxonomy above. To control for
training-data leakage we also tabulate per-case training-time accuracy
on three representative training cases, \textsc{Train-A}/\textsc{B}/\textsc{C}
(see Table~\ref{tab:main-results} for full identifiers), encountered
by every trained run because they share \texttt{split\_seed}=42.

\subsection{Experimental results}
\label{sec:exp-results}

Table~\ref{tab:main-results} reports test post-revise accuracy and
training-time accuracy on \textsc{Train-A}/\textsc{B}/\textsc{C} for
every method we trained or evaluated. Three findings stand out.

\paragraph{DPA-GRPO is the best method at both scales.}
At Qwen3-4B, DPA-GRPO reaches $0.533$ test post-revise accuracy, which
already \emph{matches} the zero-shot accuracy of the next backbone size
(Qwen3-8B zero-shot, $0.533$). At Qwen3-8B, DPA-GRPO reaches $0.572$,
the highest entry in the table; this is a $+3.9$\,pp absolute
improvement over Qwen3-8B zero-shot ($0.533$) and a $+1.9$\,pp
improvement over the strongest single-agent RL baseline (generator-only
DAPO 8B, $0.553$). The same ranking holds on every per-case
training-time metric we tabulate: at 8B, DPA-GRPO is best on
\textsc{Train-A} ($0.895$ vs.\ $0.842$), \textsc{Train-B} ($0.684$ vs.\
$0.632$), and \textsc{Train-C} ($0.579$ vs.\ $0.526$), each a
$+5{-}6$\,pp improvement over the strongest single-agent 8B baseline.
The best DPA-GRPO configuration differs by backbone (4B prefers $10$
cases per step; 8B prefers $3$); the full sweep is reported in
Table~\ref{tab:ours-ablation}, and we adopt those bolded rows as the
headline numbers throughout the paper.

\paragraph{Scaling with training data.}
Figure~\ref{fig:scaling} plots accuracy against the cumulative number of
training-line decisions ingested during training, so that runs with
different \emph{cases per step} settings are placed on a common
$x$-axis. DPA-GRPO surpasses every untrained baseline once roughly
$2{\times}10^{3}$ line decisions have been seen, and continues to
improve when given more data; generator-only GRPO (the strongest
single-agent baseline at 4B) plateaus earlier and at a lower ceiling.
The gap between DPA-GRPO and generator-only GRPO is therefore not a
data-quantity artefact of unequal training budgets, but a property of
the objective itself.

\paragraph{Scaling with backbone size.}
The two backbones tell complementary stories. At 4B, the \textsc{Train-A}
case is far from saturation ($0.526$ zero-shot), so most of the lift
DPA-GRPO obtains there ($+21.1$\,pp to $0.737$) is genuine learning.
At 8B, \textsc{Train-A} is already near-saturating zero-shot ($0.842$),
so the marginal lift on that case is small ($+5.3$\,pp to $0.895$);
the larger backbone's headroom appears instead on \textsc{Train-B} and
\textsc{Train-C}, the two harder cases, where DPA-GRPO 8B delivers
$+5.2$ and $+5.3$\,pp over the strongest single-agent 8B baseline.
This is the pattern one would hope for: as the Generator gets stronger,
DPA-GRPO's gains migrate to harder line types rather than disappearing.

\subsection{Comparisons with baselines}
\label{sec:exp-baselines}

We compare DPA-GRPO against three families of baselines, all using the
same Qwen3-4B and Qwen3-8B backbones, the same train/test split, and the
same LoRA hyperparameters. The first family is untrained; the second is
\emph{generator-only}---a single LoRA adapter trained to maximise strict
accuracy on the gold draft, with no Verifier; and the third is the
trained, two-player DPA-GRPO.

\paragraph{Untrained baselines.}
The zero-shot Generator already produces calibrated drafts on this
benchmark, with strict accuracy $0.414$ at 4B and $0.533$ at 8B
(Table~\ref{tab:main-results}, top block). Because the Verifier loop is
part of our trained method only, the zero-shot rows are intentionally
single-agent: any post-revise accuracy attributed to a zero-shot row
would conflate the loop with the underlying model. This makes
DPA-GRPO's $+11.9$\,pp lift at 4B and $+3.9$\,pp lift at 8B a
conservative measure of the parameter-update contribution.

\paragraph{Generator-only RL baselines (GRPO, DAPO, GSPO).}
We re-implemented three single-agent RL baselines that share the GRPO
trajectory-grouping recipe but differ in their objective: \emph{vanilla}
GRPO~\citep{shao2024deepseekmath} with the same clip-and-KL setup as
DPA-GRPO's Generator branch; \emph{DAPO}~\citep{yu2025dapo}, which
decouples the policy clip into asymmetric
$\varepsilon_{\text{low}}$/$\varepsilon_{\text{high}}$ bounds and adds
dynamic sampling; and \emph{GSPO}~\citep{zheng2025gspo}, which replaces
the per-token policy ratio with a per-sequence one. All three are
trained against the same gold-draft outcome reward as DPA-GRPO's
Generator branch and are matched on wall-clock budget, so the
comparison isolates the contribution of the two-player credit
assignment. At 4B, DPA-GRPO ($0.533$) beats vanilla GRPO ($0.520$),
DAPO ($0.467$), and GSPO ($0.461$). At 8B, DPA-GRPO ($0.572$) beats
vanilla GRPO ($0.539$) and DAPO ($0.553$); the GSPO 8B run is in
progress and will be reported in the camera-ready.

\paragraph{The two-player loop also stabilises the ranking.}
Beyond the absolute numbers, the ranking among single-agent baselines
\emph{flips} between scales: DAPO is the worst single-agent recipe at
4B ($0.467$) but the best at 8B ($0.553$), and GRPO does the
reverse. DPA-GRPO is the strict winner at both scales and on every
\textsc{Train-A/B/C} column we tabulate. We read this as evidence that
the two-player credit assignment is not interchangeable with a sharper
single-agent objective; it provides a complementary stabilising signal
that single-agent recipes cannot match by tuning alone.

\paragraph{Why a Verifier helps.}
A generator-only learner only ever moves probability mass between
``correct'' and ``incorrect'' draft outcomes; it has no signal that
distinguishes a \emph{confidently} correct line from a fragile one. The
two-player loop adds a Verifier whose pass-through behaviour is itself
trained, so a Generator line that is correct \emph{but} would be
overturned by an unconfident Verifier is no longer interchangeable with
a confidently-correct one in the joint reward. We make this mechanism
concrete next.

\subsection{Case discovery: which line-level dynamics change?}
\label{sec:exp-cases}

A natural worry with the headline numbers in
Table~\ref{tab:main-results} is that the post-revise gain comes from a
single, narrow mechanism: the Verifier rescuing already-wrong drafts at
test time. If that were the only thing happening, DPA-GRPO would amount
to a learned form of test-time refinement and should be compared
against post-hoc self-revision rather than against single-agent RL. To
rule this out we decompose every evaluated line into the eight-case
taxonomy of \S\ref{sec:exp-setup} and ask which buckets actually move
during training.

Figure~\ref{fig:case-improvement} reports the case-mix shift
\emph{during training}: for our two productive 4B runs, we compare the
first-5-step rolling window (baseline; the joint policy before
DPA-GRPO updates have moved either head appreciably) against the most
favourable post-update 5-step window on the same training-rollout
stream. Both endpoints share the same Verifier-firing regime, so the
comparison isolates what the joint policy \emph{learns} during training
rather than how often the Verifier intervenes. Consequently, the
percentages here are training-time within-run shifts and differ from
the test-set case mix in Table~\ref{tab:main-results} (single held-out
checkpoint); the two answer complementary questions
(\emph{what does the 2-GRPO loop learn to do?} vs.\ \emph{how does the
trained loop behave on held-out data?}). Three movements emerge, all
in the \emph{intended} direction:

\begin{itemize}
\item \textbf{Case~1 (Generator correct, Verifier silent) rises by
  $+15.1\%$.} Case~1 is the desired terminal state---the Generator
  is right and the Verifier correctly does nothing. The lift confirms
  that the Verifier's joint pressure on the Generator is producing
  more confidently-correct drafts, not just more revisions.
\item \textbf{Case~4 (\emph{missed bug}) drops by $-18.1\%$.} Case~4
  is the dominant failure mode of the initial joint policy and the only
  bucket where neither agent can be retroactively credited with a save.
  The drop here is the largest single movement we observe and is the
  mechanism by which DPA-GRPO improves \emph{strict} accuracy, not just
  post-revise accuracy.
\item \textbf{Case~6a (Generator skips a bad revision) rises by
  $+7.4\%$.} In Case~6a the draft is wrong and the Verifier correctly
  raises SAC, but the Generator's own candidate revision would also be
  wrong; the Generator chooses KEEP rather than swapping in a
  no-better revision. This is the only bucket that explicitly
  measures the Generator's \emph{discriminative} skill against its
  own revision proposals, and the rise rules out the trivial
  alternative of an always-revise Generator that would post-hoc trade
  strict accuracy for post-revise accuracy.
\end{itemize}

The three movements are mutually reinforcing rather than redundant. The
Generator becomes more right (Case~1, Case~4) and also better
calibrated against its own revision proposals (Case~6a), while the
Verifier becomes more accurate at flagging genuine bugs (the only way
Case~4 can drop
this far). The same direction of movement holds on the 8B runs
(Appendix~\ref{app:case-trends}); the within-run magnitudes are smaller
because the 8B early-training baseline already sits in a
less-Case-4-heavy regime, so there is simply less probability mass left
to redistribute during the productive training window. Taken with the test
accuracy gains in Table~\ref{tab:main-results}, this answers
question~(iv): the headline accuracy improvement is reflected in real,
predicted shifts in the line-level case mix, not in a single
narrow refinement mechanism.

% and ensure the following are also \input or \include'd before this file:
%   - paper_results_filled.tex   (Tables 1, 2 and 3)
%   - the figure block for       paper_figures/case_improvement_bars.pdf
%                                paper_figures/scaling_with_lines.pdf
%
% Required packages (NeurIPS standard, already used by the tables file):
%   \usepackage{booktabs}
%   \usepackage{multirow}
%   \usepackage{makecell}
% =============================================================================

\section{Experiments}
\label{sec:experiments}

We evaluate DPA-GRPO on TaxCalcBench TY24~\citep{bock2025taxcalcbench}, a
line-level federal-1040 reasoning benchmark, and compare against zero-shot
and trained single-agent reinforcement-learning baselines at two backbone
sizes (Qwen3-4B and Qwen3-8B). Our experiments are designed to answer
four questions, one per subsection: \emph{(i)} how is the two-player
problem framed and trained on this benchmark
(\S\ref{sec:exp-setup})?\ \emph{(ii)} how does DPA-GRPO perform overall,
and how does it scale with backbone size and training-case diversity
(\S\ref{sec:exp-results})?\ \emph{(iii)} is the gain attributable to the
two-player credit assignment, or could a stronger single-agent RL recipe
close it (\S\ref{sec:exp-baselines})?\ and \emph{(iv)} which line-level
dynamics actually change during training, and do they match the case
structure that DPA-GRPO is designed to exploit
(\S\ref{sec:exp-cases})?

\subsection{Problem setup}
\label{sec:exp-setup}

\paragraph{Benchmark.}
TaxCalcBench TY24~\citep{bock2025taxcalcbench} is a held-out collection
of U.S.\ federal tax returns whose ground truth was produced by a
deterministic tax engine. Each case provides the W-2/1099/Schedule
inputs and a 19-line gold 1040 (lines 1a, 9, 10, 11, 12, 15, 16, 19, 24,
25d, 26, 27, 28, 29, 32, 33, 34, 35a, 37). We use the official
train/test split with \texttt{split\_seed}=42 and
\texttt{test\_fraction}=0.2, yielding $n{=}8$ held-out test cases and a
total of $8\times 19 = 152$ independent line-level decisions per
evaluation. All accuracies in this section are computed against the
exact engine output for these 152 lines.

\paragraph{Two-player loop.}
Following Section~\ref{sec:method}, every line of the 1040 is treated as
a separate two-player game. The \emph{Generator} produces a draft value
for the line; the \emph{Verifier} either remains silent or raises a
\emph{Self-Assessment Critique} (SAC). When an SAC is raised, the
Generator emits a revised value and a binary \textsc{keep}/\textsc{revise}
decision. Both agents are LoRA adapters over the same frozen backbone, so
the only parameters trained by reinforcement learning are the two
adapters. The eight terminal states of the game collapse into the
mutually-exclusive case taxonomy used throughout the paper:
\textbf{Case~1} (Generator correct, Verifier silent),
\textbf{Cases~2--3} (Verifier intervenes when the original draft is
already correct),
\textbf{Case~4} (\emph{missed bug}: Generator wrong, Verifier silent),
\textbf{Cases~5a/5b} (Verifier intervenes on a wrong draft and the
revise either lands or misses gold),
and \textbf{Cases~6a/6b} (Generator overrides the Verifier; 6a is a
correct override, 6b an incorrect one).

\paragraph{Training.}
Both adapters are updated with 2-GRPO: one GRPO update on the Generator,
then one on the Verifier, repeated. Each update samples a group of
$G{=}2$ trajectories per case at temperature $0.2$ and clips the
importance ratio at $\varepsilon{=}0.2$; the shared reward is the
case-conditional shaping described in Section~\ref{sec:method}, coupled
with KL regularisation against a frozen reference
($\beta_{\text{KL}}{=}0.04$). We sweep \emph{training cases per
gradient step} $\in\{3, 5, 10\}$ for both backbones; all other
hyperparameters are held fixed (LoRA rank 16, learning rate
$1{\times}10^{-5}$, gradient checkpointing, fp16, single H100). Training
runs for at most 30 GRPO steps, with early stopping on held-out
post-revise accuracy (patience 2, minimum 10 steps).

\paragraph{Metrics.}
We report two accuracy metrics over the 152 line decisions on the
held-out test split. \textbf{Strict accuracy} is the fraction on which
the Generator's \emph{first} draft matches the gold engine output;
\textbf{post-revise accuracy} scores the Generator's \emph{final} value
after the optional revise. For zero-shot and trained single-agent
baselines the Generator is the only agent, so strict and post-revise
coincide; for DPA-GRPO we report post-revise accuracy unless otherwise
stated, since this is the metric the two-player loop optimises. We
additionally track \emph{decision accuracy} (the fraction of SAC-flagged
lines on which the Generator's keep/revise decision matches the
oracle correct-direction move) and the \emph{case-histogram} fractions
that decompose the 152 decisions into the taxonomy above. To control for
training-data leakage we also tabulate per-case training-time accuracy
on three representative training cases, \textsc{Train-A}/\textsc{B}/\textsc{C}
(see Table~\ref{tab:main-results} for full identifiers), encountered
by every trained run because they share \texttt{split\_seed}=42.

\subsection{Experimental results}
\label{sec:exp-results}

Table~\ref{tab:main-results} reports test post-revise accuracy and
training-time accuracy on \textsc{Train-A}/\textsc{B}/\textsc{C} for
every method we trained or evaluated. Three findings stand out.

\paragraph{DPA-GRPO is the best method at both scales.}
At Qwen3-4B, DPA-GRPO reaches $0.533$ test post-revise accuracy, which
already \emph{matches} the zero-shot accuracy of the next backbone size
(Qwen3-8B zero-shot, $0.533$). At Qwen3-8B, DPA-GRPO reaches $0.572$,
the highest entry in the table; this is a $+3.9$\,pp absolute
improvement over Qwen3-8B zero-shot ($0.533$) and a $+1.9$\,pp
improvement over the strongest single-agent RL baseline (generator-only
DAPO 8B, $0.553$). The same ranking holds on every per-case
training-time metric we tabulate: at 8B, DPA-GRPO is best on
\textsc{Train-A} ($0.895$ vs.\ $0.842$), \textsc{Train-B} ($0.684$ vs.\
$0.632$), and \textsc{Train-C} ($0.579$ vs.\ $0.526$), each a
$+5{-}6$\,pp improvement over the strongest single-agent 8B baseline.
The best DPA-GRPO configuration differs by backbone (4B prefers $10$
cases per step; 8B prefers $3$); the full sweep is reported in
Table~\ref{tab:ours-ablation}, and we adopt those bolded rows as the
headline numbers throughout the paper.

\paragraph{Scaling with training data.}
Figure~\ref{fig:scaling} plots accuracy against the cumulative number of
training-line decisions ingested during training, so that runs with
different \emph{cases per step} settings are placed on a common
$x$-axis. DPA-GRPO surpasses every untrained baseline once roughly
$2{\times}10^{3}$ line decisions have been seen, and continues to
improve when given more data; generator-only GRPO (the strongest
single-agent baseline at 4B) plateaus earlier and at a lower ceiling.
The gap between DPA-GRPO and generator-only GRPO is therefore not a
data-quantity artefact of unequal training budgets, but a property of
the objective itself.

\paragraph{Scaling with backbone size.}
The two backbones tell complementary stories. At 4B, the \textsc{Train-A}
case is far from saturation ($0.526$ zero-shot), so most of the lift
DPA-GRPO obtains there ($+21.1$\,pp to $0.737$) is genuine learning.
At 8B, \textsc{Train-A} is already near-saturating zero-shot ($0.842$),
so the marginal lift on that case is small ($+5.3$\,pp to $0.895$);
the larger backbone's headroom appears instead on \textsc{Train-B} and
\textsc{Train-C}, the two harder cases, where DPA-GRPO 8B delivers
$+5.2$ and $+5.3$\,pp over the strongest single-agent 8B baseline.
This is the pattern one would hope for: as the Generator gets stronger,
DPA-GRPO's gains migrate to harder line types rather than disappearing.

\subsection{Comparisons with baselines}
\label{sec:exp-baselines}

We compare DPA-GRPO against three families of baselines, all using the
same Qwen3-4B and Qwen3-8B backbones, the same train/test split, and the
same LoRA hyperparameters. The first family is untrained; the second is
\emph{generator-only}---a single LoRA adapter trained to maximise strict
accuracy on the gold draft, with no Verifier; and the third is the
trained, two-player DPA-GRPO.

\paragraph{Untrained baselines.}
The zero-shot Generator already produces calibrated drafts on this
benchmark, with strict accuracy $0.414$ at 4B and $0.533$ at 8B
(Table~\ref{tab:main-results}, top block). Because the Verifier loop is
part of our trained method only, the zero-shot rows are intentionally
single-agent: any post-revise accuracy attributed to a zero-shot row
would conflate the loop with the underlying model. This makes
DPA-GRPO's $+11.9$\,pp lift at 4B and $+3.9$\,pp lift at 8B a
conservative measure of the parameter-update contribution.

\paragraph{Generator-only RL baselines (GRPO, DAPO, GSPO).}
We re-implemented three single-agent RL baselines that share the GRPO
trajectory-grouping recipe but differ in their objective: \emph{vanilla}
GRPO~\citep{shao2024deepseekmath} with the same clip-and-KL setup as
DPA-GRPO's Generator branch; \emph{DAPO}~\citep{yu2025dapo}, which
decouples the policy clip into asymmetric
$\varepsilon_{\text{low}}$/$\varepsilon_{\text{high}}$ bounds and adds
dynamic sampling; and \emph{GSPO}~\citep{zheng2025gspo}, which replaces
the per-token policy ratio with a per-sequence one. All three are
trained against the same gold-draft outcome reward as DPA-GRPO's
Generator branch and are matched on wall-clock budget, so the
comparison isolates the contribution of the two-player credit
assignment. At 4B, DPA-GRPO ($0.533$) beats vanilla GRPO ($0.520$),
DAPO ($0.467$), and GSPO ($0.461$). At 8B, DPA-GRPO ($0.572$) beats
vanilla GRPO ($0.539$) and DAPO ($0.553$); the GSPO 8B run is in
progress and will be reported in the camera-ready.

\paragraph{The two-player loop also stabilises the ranking.}
Beyond the absolute numbers, the ranking among single-agent baselines
\emph{flips} between scales: DAPO is the worst single-agent recipe at
4B ($0.467$) but the best at 8B ($0.553$), and GRPO does the
reverse. DPA-GRPO is the strict winner at both scales and on every
\textsc{Train-A/B/C} column we tabulate. We read this as evidence that
the two-player credit assignment is not interchangeable with a sharper
single-agent objective; it provides a complementary stabilising signal
that single-agent recipes cannot match by tuning alone.

\paragraph{Why a Verifier helps.}
A generator-only learner only ever moves probability mass between
``correct'' and ``incorrect'' draft outcomes; it has no signal that
distinguishes a \emph{confidently} correct line from a fragile one. The
two-player loop adds a Verifier whose pass-through behaviour is itself
trained, so a Generator line that is correct \emph{but} would be
overturned by an unconfident Verifier is no longer interchangeable with
a confidently-correct one in the joint reward. We make this mechanism
concrete next.

\subsection{Case discovery: which line-level dynamics change?}
\label{sec:exp-cases}

A natural worry with the headline numbers in
Table~\ref{tab:main-results} is that the post-revise gain comes from a
single, narrow mechanism: the Verifier rescuing already-wrong drafts at
test time. If that were the only thing happening, DPA-GRPO would amount
to a learned form of test-time refinement and should be compared
against post-hoc self-revision rather than against single-agent RL. To
rule this out we decompose every evaluated line into the eight-case
taxonomy of \S\ref{sec:exp-setup} and ask which buckets actually move
during training.

Figure~\ref{fig:case-improvement} reports the case-mix shift
\emph{during training}: for our two productive 4B runs, we compare the
first-5-step rolling window (baseline; the joint policy before
DPA-GRPO updates have moved either head appreciably) against the most
favourable post-update 5-step window on the same training-rollout
stream. Both endpoints share the same Verifier-firing regime, so the
comparison isolates what the joint policy \emph{learns} during training
rather than how often the Verifier intervenes. Consequently, the
percentages here are training-time within-run shifts and differ from
the test-set case mix in Table~\ref{tab:main-results} (single held-out
checkpoint); the two answer complementary questions
(\emph{what does the 2-GRPO loop learn to do?} vs.\ \emph{how does the
trained loop behave on held-out data?}). Three movements emerge, all
in the \emph{intended} direction:

\begin{itemize}
\item \textbf{Case~1 (Generator correct, Verifier silent) rises by
  $+15.1\%$.} Case~1 is the desired terminal state---the Generator
  is right and the Verifier correctly does nothing. The lift confirms
  that the Verifier's joint pressure on the Generator is producing
  more confidently-correct drafts, not just more revisions.
\item \textbf{Case~4 (\emph{missed bug}) drops by $-18.1\%$.} Case~4
  is the dominant failure mode of the initial joint policy and the only
  bucket where neither agent can be retroactively credited with a save.
  The drop here is the largest single movement we observe and is the
  mechanism by which DPA-GRPO improves \emph{strict} accuracy, not just
  post-revise accuracy.
\item \textbf{Case~6a (Generator skips a bad revision) rises by
  $+7.4\%$.} In Case~6a the draft is wrong and the Verifier correctly
  raises SAC, but the Generator's own candidate revision would also be
  wrong; the Generator chooses KEEP rather than swapping in a
  no-better revision. This is the only bucket that explicitly
  measures the Generator's \emph{discriminative} skill against its
  own revision proposals, and the rise rules out the trivial
  alternative of an always-revise Generator that would post-hoc trade
  strict accuracy for post-revise accuracy.
\end{itemize}

The three movements are mutually reinforcing rather than redundant. The
Generator becomes more right (Case~1, Case~4) and also better
calibrated against its own revision proposals (Case~6a), while the
Verifier becomes more accurate at flagging genuine bugs (the only way
Case~4 can drop
this far). The same direction of movement holds on the 8B runs
(Appendix~\ref{app:case-trends}); the within-run magnitudes are smaller
because the 8B early-training baseline already sits in a
less-Case-4-heavy regime, so there is simply less probability mass left
to redistribute during the productive training window. Taken with the test
accuracy gains in Table~\ref{tab:main-results}, this answers
question~(iv): the headline accuracy improvement is reflected in real,
predicted shifts in the line-level case mix, not in a single
narrow refinement mechanism.

% and ensure the following are also \input or \include'd before this file:
%   - paper_results_filled.tex   (Tables 1, 2 and 3)
%   - the figure block for       paper_figures/case_improvement_bars.pdf
%                                paper_figures/scaling_with_lines.pdf
%
% Required packages (NeurIPS standard, already used by the tables file):
%   \usepackage{booktabs}
%   \usepackage{multirow}
%   \usepackage{makecell}
% =============================================================================

\section{Experiments}
\label{sec:experiments}

We evaluate DPA-GRPO on TaxCalcBench TY24~\citep{bock2025taxcalcbench}, a
line-level federal-1040 reasoning benchmark, and compare against zero-shot
and trained single-agent reinforcement-learning baselines at two backbone
sizes (Qwen3-4B and Qwen3-8B). Our experiments are designed to answer
four questions, one per subsection: \emph{(i)} how is the two-player
problem framed and trained on this benchmark
(\S\ref{sec:exp-setup})?\ \emph{(ii)} how does DPA-GRPO perform overall,
and how does it scale with backbone size and training-case diversity
(\S\ref{sec:exp-results})?\ \emph{(iii)} is the gain attributable to the
two-player credit assignment, or could a stronger single-agent RL recipe
close it (\S\ref{sec:exp-baselines})?\ and \emph{(iv)} which line-level
dynamics actually change during training, and do they match the case
structure that DPA-GRPO is designed to exploit
(\S\ref{sec:exp-cases})?

\subsection{Problem setup}
\label{sec:exp-setup}

\paragraph{Benchmark.}
TaxCalcBench TY24~\citep{bock2025taxcalcbench} is a held-out collection
of U.S.\ federal tax returns whose ground truth was produced by a
deterministic tax engine. Each case provides the W-2/1099/Schedule
inputs and a 19-line gold 1040 (lines 1a, 9, 10, 11, 12, 15, 16, 19, 24,
25d, 26, 27, 28, 29, 32, 33, 34, 35a, 37). We use the official
train/test split with \texttt{split\_seed}=42 and
\texttt{test\_fraction}=0.2, yielding $n{=}8$ held-out test cases and a
total of $8\times 19 = 152$ independent line-level decisions per
evaluation. All accuracies in this section are computed against the
exact engine output for these 152 lines.

\paragraph{Two-player loop.}
Following Section~\ref{sec:method}, every line of the 1040 is treated as
a separate two-player game. The \emph{Generator} produces a draft value
for the line; the \emph{Verifier} either remains silent or raises a
\emph{Self-Assessment Critique} (SAC). When an SAC is raised, the
Generator emits a revised value and a binary \textsc{keep}/\textsc{revise}
decision. Both agents are LoRA adapters over the same frozen backbone, so
the only parameters trained by reinforcement learning are the two
adapters. The eight terminal states of the game collapse into the
mutually-exclusive case taxonomy used throughout the paper:
\textbf{Case~1} (Generator correct, Verifier silent),
\textbf{Cases~2--3} (Verifier intervenes when the original draft is
already correct),
\textbf{Case~4} (\emph{missed bug}: Generator wrong, Verifier silent),
\textbf{Cases~5a/5b} (Verifier intervenes on a wrong draft and the
revise either lands or misses gold),
and \textbf{Cases~6a/6b} (Generator overrides the Verifier; 6a is a
correct override, 6b an incorrect one).

\paragraph{Training.}
Both adapters are updated with 2-GRPO: one GRPO update on the Generator,
then one on the Verifier, repeated. Each update samples a group of
$G{=}2$ trajectories per case at temperature $0.2$ and clips the
importance ratio at $\varepsilon{=}0.2$; the shared reward is the
case-conditional shaping described in Section~\ref{sec:method}, coupled
with KL regularisation against a frozen reference
($\beta_{\text{KL}}{=}0.04$). We sweep \emph{training cases per
gradient step} $\in\{3, 5, 10\}$ for both backbones; all other
hyperparameters are held fixed (LoRA rank 16, learning rate
$1{\times}10^{-5}$, gradient checkpointing, fp16, single H100). Training
runs for at most 30 GRPO steps, with early stopping on held-out
post-revise accuracy (patience 2, minimum 10 steps).

\paragraph{Metrics.}
We report two accuracy metrics over the 152 line decisions on the
held-out test split. \textbf{Strict accuracy} is the fraction on which
the Generator's \emph{first} draft matches the gold engine output;
\textbf{post-revise accuracy} scores the Generator's \emph{final} value
after the optional revise. For zero-shot and trained single-agent
baselines the Generator is the only agent, so strict and post-revise
coincide; for DPA-GRPO we report post-revise accuracy unless otherwise
stated, since this is the metric the two-player loop optimises. We
additionally track \emph{decision accuracy} (the fraction of SAC-flagged
lines on which the Generator's keep/revise decision matches the
oracle correct-direction move) and the \emph{case-histogram} fractions
that decompose the 152 decisions into the taxonomy above. To control for
training-data leakage we also tabulate per-case training-time accuracy
on three representative training cases, \textsc{Train-A}/\textsc{B}/\textsc{C}
(see Table~\ref{tab:main-results} for full identifiers), encountered
by every trained run because they share \texttt{split\_seed}=42.

\subsection{Experimental results}
\label{sec:exp-results}

Table~\ref{tab:main-results} reports test post-revise accuracy and
training-time accuracy on \textsc{Train-A}/\textsc{B}/\textsc{C} for
every method we trained or evaluated. Three findings stand out.

\paragraph{DPA-GRPO is the best method at both scales.}
At Qwen3-4B, DPA-GRPO reaches $0.533$ test post-revise accuracy, which
already \emph{matches} the zero-shot accuracy of the next backbone size
(Qwen3-8B zero-shot, $0.533$). At Qwen3-8B, DPA-GRPO reaches $0.572$,
the highest entry in the table; this is a $+3.9$\,pp absolute
improvement over Qwen3-8B zero-shot ($0.533$) and a $+1.9$\,pp
improvement over the strongest single-agent RL baseline (generator-only
DAPO 8B, $0.553$). The same ranking holds on every per-case
training-time metric we tabulate: at 8B, DPA-GRPO is best on
\textsc{Train-A} ($0.895$ vs.\ $0.842$), \textsc{Train-B} ($0.684$ vs.\
$0.632$), and \textsc{Train-C} ($0.579$ vs.\ $0.526$), each a
$+5{-}6$\,pp improvement over the strongest single-agent 8B baseline.
The best DPA-GRPO configuration differs by backbone (4B prefers $10$
cases per step; 8B prefers $3$); the full sweep is reported in
Table~\ref{tab:ours-ablation}, and we adopt those bolded rows as the
headline numbers throughout the paper.

\paragraph{Scaling with training data.}
Figure~\ref{fig:scaling} plots accuracy against the cumulative number of
training-line decisions ingested during training, so that runs with
different \emph{cases per step} settings are placed on a common
$x$-axis. DPA-GRPO surpasses every untrained baseline once roughly
$2{\times}10^{3}$ line decisions have been seen, and continues to
improve when given more data; generator-only GRPO (the strongest
single-agent baseline at 4B) plateaus earlier and at a lower ceiling.
The gap between DPA-GRPO and generator-only GRPO is therefore not a
data-quantity artefact of unequal training budgets, but a property of
the objective itself.

\paragraph{Scaling with backbone size.}
The two backbones tell complementary stories. At 4B, the \textsc{Train-A}
case is far from saturation ($0.526$ zero-shot), so most of the lift
DPA-GRPO obtains there ($+21.1$\,pp to $0.737$) is genuine learning.
At 8B, \textsc{Train-A} is already near-saturating zero-shot ($0.842$),
so the marginal lift on that case is small ($+5.3$\,pp to $0.895$);
the larger backbone's headroom appears instead on \textsc{Train-B} and
\textsc{Train-C}, the two harder cases, where DPA-GRPO 8B delivers
$+5.2$ and $+5.3$\,pp over the strongest single-agent 8B baseline.
This is the pattern one would hope for: as the Generator gets stronger,
DPA-GRPO's gains migrate to harder line types rather than disappearing.

\subsection{Comparisons with baselines}
\label{sec:exp-baselines}

We compare DPA-GRPO against three families of baselines, all using the
same Qwen3-4B and Qwen3-8B backbones, the same train/test split, and the
same LoRA hyperparameters. The first family is untrained; the second is
\emph{generator-only}---a single LoRA adapter trained to maximise strict
accuracy on the gold draft, with no Verifier; and the third is the
trained, two-player DPA-GRPO.

\paragraph{Untrained baselines.}
The zero-shot Generator already produces calibrated drafts on this
benchmark, with strict accuracy $0.414$ at 4B and $0.533$ at 8B
(Table~\ref{tab:main-results}, top block). Because the Verifier loop is
part of our trained method only, the zero-shot rows are intentionally
single-agent: any post-revise accuracy attributed to a zero-shot row
would conflate the loop with the underlying model. This makes
DPA-GRPO's $+11.9$\,pp lift at 4B and $+3.9$\,pp lift at 8B a
conservative measure of the parameter-update contribution.

\paragraph{Generator-only RL baselines (GRPO, DAPO, GSPO).}
We re-implemented three single-agent RL baselines that share the GRPO
trajectory-grouping recipe but differ in their objective: \emph{vanilla}
GRPO~\citep{shao2024deepseekmath} with the same clip-and-KL setup as
DPA-GRPO's Generator branch; \emph{DAPO}~\citep{yu2025dapo}, which
decouples the policy clip into asymmetric
$\varepsilon_{\text{low}}$/$\varepsilon_{\text{high}}$ bounds and adds
dynamic sampling; and \emph{GSPO}~\citep{zheng2025gspo}, which replaces
the per-token policy ratio with a per-sequence one. All three are
trained against the same gold-draft outcome reward as DPA-GRPO's
Generator branch and are matched on wall-clock budget, so the
comparison isolates the contribution of the two-player credit
assignment. At 4B, DPA-GRPO ($0.533$) beats vanilla GRPO ($0.520$),
DAPO ($0.467$), and GSPO ($0.461$). At 8B, DPA-GRPO ($0.572$) beats
vanilla GRPO ($0.539$) and DAPO ($0.553$); the GSPO 8B run is in
progress and will be reported in the camera-ready.

\paragraph{The two-player loop also stabilises the ranking.}
Beyond the absolute numbers, the ranking among single-agent baselines
\emph{flips} between scales: DAPO is the worst single-agent recipe at
4B ($0.467$) but the best at 8B ($0.553$), and GRPO does the
reverse. DPA-GRPO is the strict winner at both scales and on every
\textsc{Train-A/B/C} column we tabulate. We read this as evidence that
the two-player credit assignment is not interchangeable with a sharper
single-agent objective; it provides a complementary stabilising signal
that single-agent recipes cannot match by tuning alone.

\paragraph{Why a Verifier helps.}
A generator-only learner only ever moves probability mass between
``correct'' and ``incorrect'' draft outcomes; it has no signal that
distinguishes a \emph{confidently} correct line from a fragile one. The
two-player loop adds a Verifier whose pass-through behaviour is itself
trained, so a Generator line that is correct \emph{but} would be
overturned by an unconfident Verifier is no longer interchangeable with
a confidently-correct one in the joint reward. We make this mechanism
concrete next.

\subsection{Case discovery: which line-level dynamics change?}
\label{sec:exp-cases}

A natural worry with the headline numbers in
Table~\ref{tab:main-results} is that the post-revise gain comes from a
single, narrow mechanism: the Verifier rescuing already-wrong drafts at
test time. If that were the only thing happening, DPA-GRPO would amount
to a learned form of test-time refinement and should be compared
against post-hoc self-revision rather than against single-agent RL. To
rule this out we decompose every evaluated line into the eight-case
taxonomy of \S\ref{sec:exp-setup} and ask which buckets actually move
during training.

Figure~\ref{fig:case-improvement} reports the case-mix shift
\emph{during training}: for our two productive 4B runs, we compare the
first-5-step rolling window (baseline; the joint policy before
DPA-GRPO updates have moved either head appreciably) against the most
favourable post-update 5-step window on the same training-rollout
stream. Both endpoints share the same Verifier-firing regime, so the
comparison isolates what the joint policy \emph{learns} during training
rather than how often the Verifier intervenes. Consequently, the
percentages here are training-time within-run shifts and differ from
the test-set case mix in Table~\ref{tab:main-results} (single held-out
checkpoint); the two answer complementary questions
(\emph{what does the 2-GRPO loop learn to do?} vs.\ \emph{how does the
trained loop behave on held-out data?}). Three movements emerge, all
in the \emph{intended} direction:

\begin{itemize}
\item \textbf{Case~1 (Generator correct, Verifier silent) rises by
  $+15.1\%$.} Case~1 is the desired terminal state---the Generator
  is right and the Verifier correctly does nothing. The lift confirms
  that the Verifier's joint pressure on the Generator is producing
  more confidently-correct drafts, not just more revisions.
\item \textbf{Case~4 (\emph{missed bug}) drops by $-18.1\%$.} Case~4
  is the dominant failure mode of the initial joint policy and the only
  bucket where neither agent can be retroactively credited with a save.
  The drop here is the largest single movement we observe and is the
  mechanism by which DPA-GRPO improves \emph{strict} accuracy, not just
  post-revise accuracy.
\item \textbf{Case~6a (Generator skips a bad revision) rises by
  $+7.4\%$.} In Case~6a the draft is wrong and the Verifier correctly
  raises SAC, but the Generator's own candidate revision would also be
  wrong; the Generator chooses KEEP rather than swapping in a
  no-better revision. This is the only bucket that explicitly
  measures the Generator's \emph{discriminative} skill against its
  own revision proposals, and the rise rules out the trivial
  alternative of an always-revise Generator that would post-hoc trade
  strict accuracy for post-revise accuracy.
\end{itemize}

The three movements are mutually reinforcing rather than redundant. The
Generator becomes more right (Case~1, Case~4) and also better
calibrated against its own revision proposals (Case~6a), while the
Verifier becomes more accurate at flagging genuine bugs (the only way
Case~4 can drop
this far). The same direction of movement holds on the 8B runs
(Appendix~\ref{app:case-trends}); the within-run magnitudes are smaller
because the 8B early-training baseline already sits in a
less-Case-4-heavy regime, so there is simply less probability mass left
to redistribute during the productive training window. Taken with the test
accuracy gains in Table~\ref{tab:main-results}, this answers
question~(iv): the headline accuracy improvement is reflected in real,
predicted shifts in the line-level case mix, not in a single
narrow refinement mechanism.
